\begin{papercard}{sunyaev2015}
\cf{Stages} \stagetag{A} \stagetag{D} \stagetag{C}. Early large-scale study (JAMIA brief communication) establishing \emph{policy availability}, \emph{policy scope} and \emph{transparency content} as outcomes. Its scope categories are reused by Shipp and Malki.

\cf{Aim} Assess the availability, scope and transparency of mHealth app privacy policies on iOS and Android \pc{sunyaev2015}{1}{We sought to assess the availability, scope, and transparency}.

\cf{Sampling frame} The Medical and Health \& Fitness categories of both stores in May 2013. iOS data were parsed from HTML; Android data came via the open-source android-market-api because store pages were dynamically generated \pc{sunyaev2015}{1}{We selected apps from the Medical and Health and Fitness categories}. 37{,}246 apps were discovered \pc{sunyaev2015}{2}{Flow diagram for app discovery and processing}.

\cf{Relevance filter} Many category members were non-English or not health-related (e.g.\ wallpapers). 200 app descriptions were tagged manually to build a tag corpus for health functionality, information access and information handling. The remaining apps were tagged by string matching, and apps with fewer than four distinct tags were excluded \pc{sunyaev2015}{1}{We manually tagged 200 apps}. This left 24{,}405 English, health-related apps \pc{sunyaev2015}{3}{Tagging reduced the number of discovered apps to}.

\cf{Popularity rule} The 300 most frequently rated apps per store, 600 in total. Rating count was used as a proxy for downloads because iOS reports no downloads. On Android, rating count and download count correlated at Spearman $r=0.89$ \pc{sunyaev2015}{3}{since only Android (and not iOS) reports download count}.

\cf{Policy discovery} A three-step manual search, stopping at the first hit: (1) the store page, (2) the developer's website, (3) the first 30 Google results for ``\$APPNAME privacy policy'', from March to June 2013 \pc{sunyaev2015}{1}{We used a three-step manual procedure for privacy policy dis}, \pc{sunyaev2015}{3}{reviewed the first 30 results of a}.

\cf{Measures} Availability (Pearson $\chi^2$ between groups); length (word count, $t$-tests); reading grade level = mean of Flesch--Kincaid, Gunning Fog and SMOG \pc{sunyaev2015}{3}{Grade-level readability was calculated as the average of the}. \emph{Scope}, six mutually exclusive categories: single app, multiple apps, back-end application, developer homepage, all developer services, no app-related scope \pc{sunyaev2015}{3}{Privacy policy scope could be limited to the single}. \emph{Transparency content} for app-focused policies: type of information (operational, behavioural, sensitive); rationale (app operation, personalisation, secondary use); sharing (service provision, social interaction, third party); user controls (supervision, notification, correction). Operational definitions are given, e.g.\ what counts as ``sensitive'' or ``supervision'' \pc{sunyaev2015}{3}{To assess the transparency of privacy policies focusing on apps or}.

\cf{Reliability} Two researchers rated scope and content. Agreement was measured with Janson and Olsson's $\iota$, a multivariate extension of Cohen's $\kappa$ for multiple judges: $\iota = 0.94$. Remaining differences were resolved by discussion \pc{sunyaev2015}{3}{Reliability assessment with Janson}. Scope ratings were unanimous apart from naming differences \pc{sunyaev2015}{3}{Aside from initial differences in naming, privacy policy scope}.

\cf{Key findings} Only 30.5\,\% had a policy (iOS 38.3\,\% vs.\ Android 22.7\,\%). Category and price had no influence, and availability correlated only weakly with rating count \pc{sunyaev2015}{3}{Only 30.5\% of apps had privacy policies}. Mean length 1{,}755 words; reading grade level 16, i.e.\ college level \pc{sunyaev2015}{3}{Privacy policies have an average reading grade level}. 66.1\,\% did not focus on the app itself \pc{sunyaev2015}{3}{of discovered privacy policies do not focus on the app}.

\cf{Recommendations} Do not collect data not needed for the central function. Give complete transparency on disclosures and sales in a standardised format at sixth-grade reading level, with a bulleted, graphical or tabular executive summary \pc{sunyaev2015}{5}{a bulleted, graphical, or tabular executive summary should be}.

\cf{Use for this thesis} Canonical source for: the \emph{policy discovery protocol} (store, then developer site, then web search, with a stop rule); \emph{scope} categories; the argument that rating count can proxy popularity (with a reported correlation); and keeping policy \emph{absence} as an outcome with its own denominator. Readability formulas are English-specific. Rating counts measure popularity, not privacy quality.
\end{papercard}
